\begin{trapbox}
\textbf{易错点 1：原结论“面积最大值总是 $a^2$”是错误的。}
\begin{itemize}
\item 反例：$a=2,b=\sqrt3,c=1$，实际 $S_{\max}=2\sqrt3<4$。
\item 正确分类：$c\le b$ 时 $2bc$；$c\ge b$ 时 $a^2$。
\item 仅有 $S\le a^2$ 是上界，\textbf{还必须检验取等条件}。
\end{itemize}

\textbf{易错点 2：截距积并非在所有切点都有定义。}
\begin{itemize}
\item $x_0=0$（上下顶点）：切线平行于 $x$ 轴，$N=l\cap Ox$ 不存在。
\item $y_0=0$（左右顶点）：$F_1,P,F_2$ 共线，焦点角退化；法线与 $x$ 轴重合，``法线交 $x$ 轴于唯一点 $M$'' 失效。
\item 故 $x_Mx_N=c^2$ 的严格使用条件是 $x_0y_0\ne0$。
\end{itemize}

\textbf{易错点 3：把角平分线记反。}

非长轴顶点处，\textbf{法线}平分 $\angle F_1PF_2$ 的内角；
\textbf{切线}平分外角。光线入射、反射时，要按传播方向辨别入射线与反射线，
不要直接把 $PF_1$ 和 $PF_2$ 两条射线误看成传播方向相同。

\textbf{易错点 4：反射证明出现循环论证。}

不能直接声称 ``$F_1$ 关于切线的像 $Q$ 与 $F_2,P$ 共线''，
再把这句话当作光学性质的证明；\textbf{这本来就是待证性质}。
需要先用椭圆的最短路/三角不等式，或等值曲线梯度，证明共线性。

\textbf{易错点 5：把 ``延长 $F_1P$，取等长点'' 与 ``沿切线反射'' 混为一谈。}

正确的 $Q$ 是 $F_1$ 关于\textbf{切线 $l$ 的镜像}，满足 $PF_1=PQ$、
$H_1$ 是 $F_1Q$ 中点；不是随便在 $F_1P$ 延长线上取一个点。
否则一般不能推出 $F_2Q=2a$，中位线证明也不能成立。

\textbf{易错点 6：焦点到切线的距离公式漏绝对值或把它们相减。}

对 $Ax+By=1$，
\[
d_1=\frac{1+cA}{\sqrt{A^2+B^2}},\quad
 d_2=\frac{1-cA}{\sqrt{A^2+B^2}}
\]
是因为 $|cA|<1$；两个焦点位于切线同侧，距离都是正值。
不能无依据地把 $d_1+d_2$ 改成差。

\textbf{易错点 7：斜率法遗漏竖直切线。}

$y=kx+m$ 只覆盖非竖直切线。
端点 $P=(\pm a,0)$ 对应 $x=\pm a$，应改用
$\frac{x_0x}{a^2}+\frac{y_0y}{b^2}=1$ 或单独验算。

\textbf{易错点 8：几种“圆”并非同一个概念。}

$H_1,H_2$ 所在的是\textbf{以椭圆中心为圆心、半径为 $a$ 的辅助圆}。
它不是以焦点为圆心、半径为 $b$ 的圆；
也不意味着切点 $P$ 在该圆上（除长轴顶点外，通常不在）。

\textbf{易错点 9：四边形 $F_1H_1H_2F_2$ 的面积公式不能乱用。}

相互平行的是 $F_1H_1$ 与 $F_2H_2$（两条焦点垂线），
并非一般情况下的 $F_1F_2$ 与 $H_1H_2$。
故梯形面积为 $S=\frac12(d_1+d_2)H_1H_2$，
竖直切线时两垂足重合且面积为 $0$。

\textbf{易错点 10：点幂定理需要明确第二个交点。}

不能误认 $H_1,F_1,H_2$ 共线；真正与 $H_1$ 同弦的是
$K=-H_2$（$H_2$ 的\textbf{对径点}），所以 $F_1K=d_2$。
此时点幂的正长度形式是 $F_1H_1\cdot F_1K=a^2-c^2$。

\textbf{易错点 11：推极限 $c\to0$ 时要区分公式与图形。}

圆的情形两焦点重合，$d_1d_2=a^2$ 仍有自然极限，
但焦点角退化，角平分线与两个截距的几何配置不再按原方式定义。

\textbf{易错点 12：TeX 书写规范。}

带下标的符号必须放入数学环境，如 $d_1,d_2,x_M,x_N$；
不要把美元符号写成普通文字。面积分段公式应放在数学环境中。
\end{trapbox}
